% from main.tex
%%%%%%%%%%  scan idx 558  |  folio 545  %%%%%%%%%%
\end{equation}
Rappelons la définition du degré \(n\) de \(f\) [3] : il est nul sauf dans le cas où \(\dim(X)=\dim(Y)=\dim(f(Y))\). Dans ce dernier cas on pose
\begin{equation}\tag{4.4.4}
 n=(\len_{\mathcal O_{X,x}}(\mathcal O_{Y,y}))/(\len_{\mathcal O_{X,x}}(\mathcal O_{X,x})),
\end{equation}
où \(x\) est le point générique de \(X\) et \(y\) le point générique de \(Y\).
\end{enumerate}

Dans la démonstration de (4.4.3) le cas \(\dim(X)=\dim(Y)\) est le seul non-trivial. Si \(\dim(f(Y))<d\), alors on a
\begin{equation}\tag{4.4.5}
 R^if_*(\mathcal O_Y)=0
\end{equation}
pour \(i>0\) dans un voisinage de \(x\). Mais la condition (4.4.5) est satisfaite aussi dans le cas \(\dim(f(Y))=d\), car alors la dimension de la fibre générique est nulle (EGA III 4.2.2 et EGA IV 5.6.6). Alors 1.1.2 donne les congruences
\[
 f_*(\cl_\bullet(\mathcal O_Y))
 \equiv \cl_\bullet(f_*(\mathcal O_Y)
 \equiv
 \len_{\mathcal O_{X,x}}(\mathcal O_{Y,y})
 \cl_\bullet(\mathcal O_{X_{\red}})
\]
modulo \(\Filt_{d-1}(X)\). En vertu de (4.4.2) il en découle que
\[
 (f^*L_1,\ldots,f^*L_d;\mathcal O_Y)
 =
 \len_{\mathcal O_{X,x}}(\mathcal O_{Y,y})
 (L_1,\ldots,L_d;\mathcal O_{X_{\red}}).
\]
D'autre part \(2^\circ\) donne
\[
 (L_1,\ldots,L_d;\mathcal O_X)
 =
 \len_{\mathcal O_{X,x}}(\mathcal O_{X,x})
 (L_1,\ldots,L_d;\mathcal O_{X_{\red}}),
\]
d'où l'assertion.

\subsection*{4.5}
On
